\subsection{Generation, Retrieval and Feedback}
The task-level landscape in Table~\ref{tab:comparison_landscape} distinguishes output, supervision and feedback rather than ranking performance. OptiMUS constructs formulations/code with feedback, while PaMOP partitions the description and models grouped constraints with augmented prompts and iterative correction \cite{ahmaditeshnizi2024optimus_icml,pan2025pamop}. OptMATH generates data by back-translation and forward checking; ORLM trains modeling language models and DeepOR adds intermediate reasoning and solver-feedback learning \cite{lu2025optmath,huang2025orlm,xiao2026deepor}. Their outputs are richer than supplied-slot mappings.

OPT2CODE uses retrieval for code generation \cite{ahmed2025opt2code}. Autoformulation, uncertainty-aware modeling and Chain-of-Experts address formulation decisions or coordinated agents \cite{astorga2025autoformulation,yu2026uncertainty,xiao2024chainofexperts}. OptLLM and LLMOPT combine modeling with code/feedback \cite{zhang2024optllm,jiang2025llmopt}; SIRL and OR-R1 use solver-informed learning \cite{chen2025solver,ding2026orr1}. Dual-cluster, experience-replay and solver-independent approaches concern retained knowledge, correction or representation \cite{zhang2026dual,xie2026escaping,li2026solver}. AlphaOPT and ORThought extend the context through accumulated experience or structured reasoning \cite{kong2025alphaopt,yang2026orthought}. OptimAI (Thind et al., unpublished preprint, \href{https://arxiv.org/abs/2504.16918}{arXiv:2504.16918}) studies multi-agent formulation, planning and code revision. None is assigned an unexecuted accuracy result here.

\subsection{Other Grounding and Validation Context}
Spider and the survey by Affolter et al. contextualize database interfaces \cite{yu2018spider,affolter2019nli}. ReMatch (Sheetrit et al., unpublished preprint, \href{https://arxiv.org/abs/2403.01567}{arXiv:2403.01567}) uses retrieval-enhanced schema matching. Execution-guided decoding (Wang et al., unpublished preprint, \href{https://arxiv.org/abs/1807.03100}{arXiv:1807.03100}) uses partial Structured Query Language (SQL) execution to reject some faulty programs. MIPLIB-NL addresses more complex optimization formulations \cite{miplibnl2026}. These settings provide methodological context, not interchangeable scalar-grounding scores.
